% 03-core.tex: ядро удостоверения.
\section{Ядро: одно удостоверение, одна подпись, один человек}
\label{sec:core}

\motif{Приватность не есть свобода от принуждения.}

\subsection{Что такое удостоверение}

В центре Полярис находится одна строка таблицы \entity{ТокенЛичности} и по меньшей мере одна
строка \entity{ПодписьТокена}. Строка токена записывает уникальное \code{значение\_токена},
физический серийный номер, человека, которому он принадлежит, ведомство, которое его выдало,
алгоритм, под которым он был выдан, его предшественника, если он заменил собой более ранний
токен, и его состояние. Строка подписи записывает подпись МЛ-ДСА выдавшего ведомства над дайджестом
ША3-256 значения токена вместе с открытым ключом, которым она была произведена, так что проверка
самодостаточна и переживает вывод подписывающего ключа из обращения. Токен без строки подписи
существовать не может: триггер базы данных отвергает любую запись, которая оставила бы токен без
неустаревшей подписи. Удостоверение есть поэтому не документ, а подписанное утверждение:
\emph{этот орган выдал этот токен этому человеку.} Утверждение может также
связывать ключ, которым распоряжается держатель, записанный в \entity{СобытиеКлючаДержателя} под той же
дисциплиной добавления без изменения, что и ключ удостоверяющего органа, так что владение файлом
перестаёт быть владением удостоверением (раздел~\ref{sec:agents}). \sImpl

Держатель держит у себя пакет подлинности: значение токена, подпись и открытый ключ эмитента,
выгружаемые приложением и хранимые в виде файла эталонным кошельком
(\code{сценарии/\allowbreak{}полярис-кошелёк.пи}). Любой может проверить этот пакет стандартной библиотекой
МЛ-ДСА и без единой строки кода Полярис (раздел~\ref{sec:verifiers}).

\subsection{Физический токен как предмет, который способен изготовить кто-то другой}

Кристалл смоделирован и не произведён; всё, что его окружает, существует. Карта есть кодировка с эталонным кодеком и опубликованными векторами
(\code{полярис\_карта/\allowbreak{}профиль\_карты.пи}), программный токен, говорящий на ИСО 7816-4 настоящими
словами состояния, с двумя гнёздами для ПИН, счётчиком попыток и ПУК (\code{эмулятор.пи}), шаг
персонализации, на котором карта сама вырабатывает свои ключи, так что орган никогда ими не
владеет, и эталонное устройство проверки, которое читает карту по НФС либо предъявление по
куар-коду и сообщает три вывода по отдельности: владение, подлинность и полномочия, потому что
отозванная карта всё равно подпишет вызов, а устройство, говорящее лишь \emph{принято}, ничему
своего оператора не учит. Профиль прямо называет своё единственное формообразующее ограничение:
сертифицированные защищённые элементы, подписывающие по ФИПС 204 аппаратно, сегодня не есть
нечто, что орган может закупить массово, поэтому карта несёт две подписи эмитента над одним телом,
классическую, которую нынешний кремний произвести способен, и МЛ-ДСА, а считыватель решает, какой
из них доверяет. Каждая из этих вещей выполняется в учении при каждой отправке. \sImpl{} для
профиля, эмулятора, векторов и устройства; \sFut{} для произведённой карты, чего никакое учение
установить не может.

\figfile{nucleus}{Геометрия сущностей ядра, не изменившаяся с версии 1. Семь сущностей вокруг
\entity{ТокенЛичности}: человек, который его держит, ведомство, которое его выдаёт, алгоритм,
который его подписывает, контексты, в которых он может проверяться, и два журнала событий,
работающих только на добавление, которые он порождает. Связи <<многие ко многим>> разрешаются в
соединительные таблицы \entity{ДопускКАлгоритму} и \entity{РазрешениеТокена}. Эта геометрия есть
постоянная величина документа: раздел~\ref{sec:walls} рисует остальные тридцать восемь таблиц как
кольца вокруг неё.}{fig:nucleus}

\subsection{Чем удостоверение не является}

Десятое ограничение конституции, С10, гласит, что удостоверение личности никогда не несёт
платёжных полномочий: в схеме нет ни одного денежного поля, и инвариантная проверка роняет
сборку, если такое появится. Причина изложена как режим отказа, а не как предпочтение.
Удостоверение, которое есть ещё и деньги, наследует программируемость: всякое ограничение, какое
кому-либо хочется наложить на траты, нарастает на идентификационный токен, покуда
административная ошибка в бумагах не станет экзистенциальной ошибкой в банковском балансе и покуда
системе нельзя будет сказать, что этому человеку нельзя покупать топливо. Тот же отказ покрывает
репутацию, поведение и профили. Полярис отвечает на один вопрос, \emph{действителен ли этот токен
для этого контекста прямо сейчас}, и уклоняется от ответа на вопрос \emph{следует ли позволить
этому человеку сделать то или иное}. Второй вопрос принадлежит доверяющей стороне, за границей. \sImpl

\subsection{Один действующий токен на человека}

Человек вправе держать несколько подготовленных токенов, но лишь один может быть \code{ДЕЙСТВУЮЩИЙ}.
Два действующих токена у одного человека позволили бы одному токену разрешить операцию, а другому
от неё отпереться. Это правило С3, и оно закреплено частичным уникальным индексом
\code{ун\_один\_действующий\_на\_человека} по идентификатору человека там, где состояние есть
действующее. Два оператора, активирующие два резервных токена для одного держателя в один и тот
же миг, обнаруживают, что удаётся ровно одному из них, и тест, который это доказывает, работает
настоящими потоками против живой базы данных (С9). Правило также смоделировано в ТЛА+, и модель
проверяется при каждой отправке, при всяком чередовании одновременных выдачи и отзыва, вместе с
парной конфигурацией, на которой она обязана упасть. Верно про него и ещё одно, записанное в
сопутствующем \emph{Протоколе свидетельств}: когда индекс убрали, чтобы посмотреть, кто это
заметит, выяснилось, что набор тестов, существующий ради доказательства С3, сам его и нарушал,
фиксируя второй действующий токен перед тем, как упасть. \sImpl

То, что С3 называет лицом, есть строка, \code{ФизическоеЛицо.физ\_лицо\_ид}, а не человек. Индекс делает
невозможным второе действующее удостоверение у той же строки; он ничего не делает с одним и тем же
человеком, зачисленным двумя строками, и никакой механизм в репозитории этого не делает.
Уникальность людей держится на зачислении, где оператор записывает свидетельства, а уровень
достоверности выводится из них.

\subsection{Жизненный цикл}

\figfile{lifecycle}{Автомат жизненного цикла токена в том виде, в каком его закрепляет триггер
\code{тр\_автомат\_состояний\_токена}: шесть законных переходов, каждый подписан событием жизненного
цикла, которое он автоматически записывает. \code{ДЕЙСТВУЮЩИЙ} есть единственное рабочее состояние.
Токен входит через \code{РЕЗЕРВ} и активируется либо отзывается, так и не будучи активирован;
действующий токен уходит через одно из трёх конечных состояний либо в \code{СПЯЩИЙ}, когда дело
принимает преемник.}{fig:lifecycle}

Триггер закрепляет переходы, а второй триггер записывает строку \entity{СобытиеЦиклаТокена} на каждую смену состояния, включая
смену, сделанную вручную из консоли базы данных. Преемство идёт по ссылке, а не перезаписью:
замещающий токен указывает на своего предшественника, а старый токен остаётся в базе данных в
своём конечном состоянии. Утраченные токены остаются утраченными, и их данные не стираются.
\sImpl

\subsection{Вход, выход и возвращение}

Три механизма окружают ядро, чтобы схему нельзя было обратить против держателя на краях
жизненного цикла, и каждый хранит запись о том, на чём он держался.

\begin{description}
  \item[Вход: многоуровневая регистрация и то, на чём держится запись.] У каждого человека есть
    состояние регистрации из пяти значений, и \code{ОСВОБОЖДЁН} есть полноправное из них: признанный
    способ участвовать в гражданской жизни без токена, будь то по биометрической несовместимости,
    религиозному освобождению или по иному пути. Запись об освобождении есть одна строка;
    перечислить незарегистрированных возможно, но ни одна функция в этом не помогает, а всякое
    такое перечисление заносится в аудит. Удостоверение записывает, как человек был подтверждён
    как тот, кем себя называет: \entity{ПодтверждениеРегистрации} записывает одно событие подтверждения личности, а
    \entity{ОснованияРегистрации} то, на чём оно держалось, обе только с добавлением записей, и
    уровень достоверности \emph{выводится} из свидетельств, а не вбивается оператором; поля для
    самого документа намеренно нет, а названный набор полей отвергается наотрез, так что запись не
    способна воссоздать паспорт. Чего это не устанавливает, изложено в пакете для рецензента как
    ограничение: подлинность свидетельств принимается со слов оператора. \sImpl

    Всякое сочетание свидетельств, какое называет любой стандарт, начинается с документов, и
    человек, у которого их нет, выпускник детского дома, беженец, взрослый, никогда не имевший
    паспорта, не проходит все эти сочетания по умолчанию. Доверенное поручительство
    (\entity{Поручительство}) есть ответ на это, и оно же есть простейший способ вычеканить
    уровень достоверности из ничего, поэтому репозиторий обращается с механизмом против
    исключения и каналом подлога как с одной и той же таблицей: поручительство ограничено,
    совместно подписывается сверх порога и невидимо на самом удостоверении, потому что человек,
    которому оно понадобилось, не должен носить за это метку у каждой стойки. Регистрационный
    код, последняя часть, не требующая никакого железа, устанавливает лишь то, что некто,
    имеющий доступ к каналу, вернул секрет, и уровень, который он способен внести, ограничен
    соответственно. \sImpl
  \item[Выход: усмотрение эмитента как решение, переживающее собственное изменение.] Худший отказ
    идентификационной системы есть не поддельный токен, а орган, отзывающий токены у целого
    населения в промышленном масштабе. Триггер ограничивает скорость, с которой одно ведомство
    способно отзывать собственные токены (по умолчанию пять процентов действующей базы за
    тридцать дней), а сверх того требуется соподписант из другого ведомства, заносимый в след
    аудита. Пять процентов есть значение развёртывания по умолчанию и параметр управления, а не
    инвариант безопасности: орган вправе задать своё (одно из ведомств образца работает при семи), и
    никакое свойство системы из самого числа не следует. Ограничение превращает неожиданность в тенденцию, которую отчётность видит; оно не
    останавливает медленного злоупотребления под потолком, о чём и говорит. Настройки вокруг
    этой власти суть записи: сам потолок (\entity{ПолитикаУсмотрения})
    можно изменить только с изложенным обоснованием, и он оставляет строку о том, кто его изменил
    и с какого значения; квота на то, какую долю населения ведомство вправе затронуть, есть
    решение с причиной, а не величина, которую можно перезаписать; строка в корне иерархии,
    \entity{Ведомство}, несёт триггер, записывающий всякое изменение уровня органа и отмечающий, было
    ли оно расширением; а планка доверяющей стороны фиксируется с той же дисциплиной, с пометкой об ослаблении и с
    обязательными причинами. \sImpl
  \item[Возвращение: обряд восстановления.] Пожар забирает кошелёк вместе с резервом.
    Восстановление обязательно и вместе с тем есть путь выдать себя за другого, поэтому его делают
    дорогим, не делая невозможным: двухфазный обряд с сорокавосьмичасовым периодом выдержки, три
    независимых внеполосных канала, утверждающий, который не может быть заявителем, и завершение
    только администратором. Поля решения движутся в одну сторону под триггером, а правило
    четырёх глаз, период выдержки и три канала суть отказы хранимой процедуры, и мутационное учение
    показывает про каждый, что тесты заметили бы его утрату (раздел~\ref{sec:walls}). Ни один путь в
    продукте не записывает три канала (раздел~\ref{sec:limitations}). Держатель остаётся без личности на время выдержки; замещающее
    удостоверение, которого промышленная система на это окно захотела бы, не построено. \sImpl{}
    для обряда, \sFut{} для временного удостоверения.
\end{description}

\subsection{Принуждение: сигнал, которого принуждающий не видит}

Призвание, стоящее над десятью ограничениями, есть противодействие принуждению: никого нельзя
заставить отречься от своей личности, передать её или сдать против собственной воли. Механизм,
дающий этому предложению опору, стар и пришёл из банковского дела: второй секрет, выглядящий для
всякого наблюдающего ровно как успех, в то время как тревога уходит по каналу, которого
наблюдающий не видит. Держатель вправе завести код принуждения. Когда он предъявлен, проверка идёт
одинаково на каждой видимой оператору поверхности, та же страница, то же всплывающее сообщение, та
же строка события, а строка \entity{СобытиеПринуждения} добавляется в таблицу, к которой экраны оператора
никогда не присоединяются и из которой пейджер срабатывает с наивысшей важностью. Сравнение
выполняется за постоянное время, и работа оплачивается на обоих путях. На карте тот же механизм
есть второй ПИН той же длины, что и первый, вводимый тем же способом и отпирающий второе гнездо
ключа, открытого ключа которого объект карты не несёт, а учение измеряет отклик и во времени, и в
байтах, потому что принуждающий у считывателя наблюдает, сколько карта провозилась. \sImpl{} как
проверенное свойство смоделированного потока.

Механизм принуждения \emph{осведомлён о принуждении}, а не устойчив к нему: он способен скрыть
тревогу от принуждающего, не знающего о его существовании, и не защищает ни от осведомлённого
принуждающего, ни от институционального доступа. Проверка отвергает более сильное слово на каждой
внешней поверхности. Против беспечного принуждающего, который не знает о существовании кодов
принуждения, механизм работает: экран неотличим, а тревога беззвучна. Против осведомлённого
принуждающего, стоящего рядом с держателем, он не работает: код набирается на терминале
верификатора, нажатия клавиш не происходят за постоянное время, а поскольку заведение кода
добровольно, \emph{у меня нет кода принуждения} есть утверждение, на которое принуждающий может
надавить и которого держатель доказать не может. Против последующего законного или ведомственного
доступа он хуже, чем ничего: \entity{СобытиеПринуждения} работает только на добавление, без пути к
очистке, что есть С1, работающее как задумано, и потому там, где принуждающая сторона есть само
учреждение или способна позже его принудить, механизм изготавливает неизгладимое свидетельство
того, что держатель сопротивлялся. Оценка не предлагает изменения механизма, и настоящий документ
не предлагает тоже; он фиксирует границу. Утверждение о побочных каналах, что ничто наблюдаемое не
отличает два пути, измерено репозиторием и никем более. \sExt

\begin{layercard}{Карточка слоя: ядро удостоверения}
\qa{Что это}{Одна строка \entity{ТокенЛичности} и её строки \entity{ПодписьТокена}, необязательно связанные с ключом держателя: подписанное утверждение, что некий орган выдал этот токен этому человеку, в жизненном цикле с одним действующим состоянием, с записанным рядом подтверждением личности, на котором оно держалось.}
\qa{Зачем оно существует}{Всякая иная структура в Полярис устраивается вокруг этого утверждения. Оно держится настолько малым, насколько возможно, чтобы остальное можно было ограничить.}
\qa{Чему доверяет}{Подписывающему ключу выдавшего ведомства, хранимому за интерфейсом хранения (раздел~\ref{sec:signatures}); ограничениям схемы; слову оператора о том, что предъявленные свидетельства были подлинными.}
\qa{Чему отказывает}{Коду приложения как источнику истины: автомат состояний, уникальность и аудит решаются базой данных. Уровню достоверности, вбитому, а не выведенному. Деньгам, репутации и профилям: С10.}
\qa{Что видит}{Законное имя держателя, дату рождения и страну; выдавшее ведомство; алгоритм; жизненный цикл; род свидетельств, на которых держалась регистрация.}
\qa{Чего не видит}{Биометрический шаблон (никогда не хранится), документ, стоящий за регистрацией (поля для него не существует), граф проверок (раздел~\ref{sec:privacy}), историю трат (никогда не моделировалась).}
\qa{Если откажет}{Поддельное ядро есть поддельная подпись, которую отвергают проверка двумя свидетелями и отделяемый верификатор; сдвоенное ядро отвергается уникальным индексом; переписанное ядро отвергается триггерами добавления; принуждающий, который сам есть учреждение, не останавливается, и запись об этом говорит.}
\qa{Кем проверяется}{\code{проверка\_объекты\_с1с10\_разрешаются}, тесты свойств С3 и модель С3, набор тестов на одновременность, \code{проверка\_профиль\_карты}, \code{проверка\_эмулятор\_карты}, \code{проверка\_устройство\_проверки}, \code{проверка\_принуждение\_на\_карте}, \code{проверка\_утверждения\_о\_принуждении\_осведомлённые}, случаи подлога и подмены в наборе атак.}
\qa{Сейчас или потом}{Сейчас. Это та часть Полярис, которая существует с версии 1; карта есть предмет уже сейчас, а кремний потом.}
\qa{Внутри и снаружи}{Внутри: человек, существующий за пределами программ. Снаружи: стена ограничений (раздел~\ref{sec:walls}).}
\end{layercard}

\analogy{Городская ратуша держит один реестр. Каждая запись говорит, кто этот человек, какой
чиновник её внёс, под какой печатью и какие бумаги чиновнику предъявили, но никогда не хранит
самих бумаг. Реестр нельзя править, к нему можно только дописывать; вторая запись о том же
человеке отвергается прямо у стойки; а печать такова, что всякий читатель вправе сверить её с
опубликованным знаком чиновника, ни о чём ратушу не спрашивая.}
